% from ega1-4-fr.tex
\oldpage[I]{124}
  que sa restriction à chacun des préschémas induits $f^{-1}(V_\lambda)$ soit
  une immersion (resp. une immersion ouverte) dans $V_\lambda$.
  \item[b)] Soient $f$ un morphisme $Y\to X$, $(V_\lambda)$ un recouvrement
  ouvert de $X$. Pour que $f$ soit une immersion fermée, il faut et il suffit
  que sa restriction à chacun des préschémas induits $f^{-1}(V_\lambda)$ soit
  une immersion fermée dans $V_\lambda$.
\end{enumerate}
\end{corollary}

Soit $f=(\psi,\theta)$~; dans le cas a), $\theta_y^\sharp$ est surjectif
(resp. bijectif) pour tout $y\in Y$, et dans le cas b) $\theta_y^\sharp$ est
surjectif pour tout $y\in Y$~; il suffit donc de vérifier que $\psi$, dans le
cas a), est un homéomorphisme de $Y$ sur une partie localement fermée (resp.
ouverte) de $X$, et dans le cas b), un homéomorphisme de $Y$ sur une partie
fermée de $X$. Or, $\psi$ est évidemment injectif et transforme tout voisinage
de $y$ dans $Y$ en un voisinage de $\psi(y)$ dans $\psi(Y)$ pour tout $y\in Y$,
en vertu de l'hypothèse~; dans le cas a), $\psi(Y)\cap V_\lambda$ est localement
fermé (resp. ouvert) dans $V_\lambda$, donc $\psi(Y)$ est localement fermé
(resp. ouvert) dans la réunion des $V_\lambda$, et \emph{a fortiori} dans $X$~;
dans le cas b), $\psi(Y)\cap V_\lambda$ est fermé dans $V_\lambda$, donc
$\psi(Y)$ est fermé dans $X$ puisque $X=\bigcup_\lambda V_\lambda$.

\begin{proposition}[4.2.5]
\label{I.4.2.5-fr}
Le composé de deux immersions (resp. de deux immersions ouvertes, de deux
immersions fermées) est une immersion (resp. une immersion ouverte, une
immersion fermée).
\end{proposition}

Cela résulte trivialement de (\hyperref[I.4.1.6-fr]{4.1.6}).

\subsection{Produit d'immersions.}
\label{subsection:I.4.3-fr}

\begin{proposition}[4.3.1]
\label{I.4.3.1-fr}
Soient $\alpha:X'\to X$, $\beta:Y'\to Y$ deux $S$-morphismes~; si $\alpha$ et
$\beta$ sont des immersions (resp. des immersions ouvertes, des immersions
fermées), $\alpha\times_S\beta$ est une immersion (resp. une immersion ouverte,
une immersion fermée). En outre, si $\alpha$ (resp. $\beta$) identifie $X'$
(resp. $Y'$) à un sous-préschéma $X''$ (resp. $Y''$) de $X$ (resp. $Y$),
$\alpha\times_S\beta$ identifie l'espace sous-jacent à $X'\times_S Y'$ au
sous-espace $p^{-1}(X'')\cap q^{-1}(Y'')$ de l'espace sous-jacent à
$X\times_S Y$, en désignant par $p$ et $q$ les projections de $X\times_S Y$
dans $X$ et $Y$ respectivement.
\end{proposition}

Compte tenu de la déf. (\hyperref[I.4.2.1-fr]{4.2.1}), on peut se restreindre
au cas où $X'$ et $Y'$ sont des sous-préschémas, $\alpha$ et $\beta$ les
morphismes d'injection. La proposition a déjà été établie pour les
sous-préschémas induits sur les ouverts (\hyperref[I.3.2.7-fr]{3.2.7})~; comme
tout sous-préschéma est un sous-préschéma fermé d'un préschéma induit sur un
ouvert (\hyperref[I.4.1.3-fr]{4.1.3}), on est ramené au cas où $X'$ et $Y'$ sont
des sous-préschémas \emph{fermés}.

Montrons d'abord qu'on peut supposer que $S$ est \emph{affine}. En effet, soit
$(S_\lambda)$ un recouvrement de $S$ formé d'ouverts affines~; si $\varphi$ et
$\psi$ sont les morphismes structuraux de $X$ et $Y$, soit
$X_\lambda=\varphi^{-1}(S_\lambda)$, $Y_\lambda=\psi^{-1}(S_\lambda)$. La
restriction $X'_\lambda$ (resp. $Y'_\lambda$) de $X'$ (resp. $Y'$) à
$X_\lambda\cap X'$ (resp. $Y_\lambda\cap Y'$) est un sous-préschéma fermé de
$X_\lambda$ (resp. $Y_\lambda$), les préschémas $X_\lambda$, $Y_\lambda$,
$X'_\lambda$, $Y'_\lambda$ peuvent être considérés comme des
$S_\lambda$-préschémas et les produits $X_\lambda\times_S Y_\lambda$ et
$X_\lambda\times_{S_\lambda}Y_\lambda$ (resp.
$X'_\lambda\times_S Y'_\lambda$ et
$X'_\lambda\times_{S_\lambda}Y'_\lambda$) sont identiques
(\hyperref[I.3.2.5-fr]{3.2.5}). Si la proposition est vraie lorsque $S$ est
affine, la restriction de $\alpha\times_S\beta$ à chacun des
$X'_\lambda\times_S Y'_\lambda$ sera donc une immersion
(\hyperref[I.3.2.7-fr]{3.2.7}). Comme le produit
$X'_\lambda\times_S Y'_\mu$ (resp. $X_\lambda\times_S Y_\mu$) s'identifie à
$(X'_\lambda\cap X'_\mu)\times_S(Y'_\lambda\cap Y'_\mu)$ (resp.
$(X_\lambda\cap X_\mu)\times_S(Y_\lambda\cap Y_\mu)$)
(\hyperref[I.3.2.6.4-fr]{3.2.6.4}), la restriction de
$\alpha\times_S\beta$
